%!TEX root = ../dokumentation.tex

\chapter{Einleitung}

Beschlagene Windschutzscheiben stellen ein Sicherheitsrisiko dar und erhöhen gleichzeitig 
den Energiebedarf von Fahrzeugen. Besonders in Elektrofahrzeugen kann der Einsatz des 
Heiz-, Lüftungs- und Klimasystems \ac{HVAC} System die Reichweite um bis zu 13\% reduzieren 
\citeM{i}{SenReich.2024}{}{}. Da die Enteisung oder Entfeuchtung der Scheibe meist 
reaktiv erfolgt, entsteht kurzfristig ein hoher Energieverbrauch.

Ein präventiver Ansatz könnte diesen Bedarf reduzieren: Wenn Temperatur- und 
Feuchtigkeitssensoren frühzeitig erkennen, dass die Bedingungen für eine Kondensation 
erreicht werden, lassen sich geeignete Maßnahmen bereits vor dem Beschlag einleiten.

Ziel dieser Arbeit ist es daher, das Potenzial eines auf ESP32 und DHT11 basierenden 
Sensorsystems zur frühzeitigen Erkennung von Beschlagbedingungen zu untersuchen und 
dessen Eignung im Hinblick auf Energieeffizienz zu bewerten.

